\documentclass[a4paper]{article}

\usepackage[UTF8]{ctex}
\usepackage[a4paper,margin=1in]{geometry}
\usepackage{graphicx}
\usepackage{float}
\usepackage{listings}
\usepackage{longtable}
\usepackage{booktabs}
\usepackage{hyperref}
\usepackage{fancyhdr}
\usepackage{lastpage}
\usepackage{color}
\usepackage{indentfirst}        % 首段缩进
\setlength{\parindent}{2em}     % 缩进2字符
\usepackage{zhnumber}           % 中文编号
\usepackage[dvipsnames]{xcolor}

%\newcommand{\college}{中山大学计算机学院}
\newcommand{\projname}{操作系统原理实验| Lab0}
\newcommand{\reporttitle}{环境搭建与实验准备}
\newcommand{\stuno}{23336188}
\newcommand{\authorname}{欧阳易芃}
\newcommand{\major}{计算机科学与技术}
\newcommand{\adviser}{陈鹏飞}
\newcommand{\startdate}{2025年3月1日}
\newcommand{\labenddate}{2025年3月2日}
\newcommand{\labroom}{超算中心307}

\pagestyle{fancy} % 使用 fancyhdr 风格
\fancyhf{}      % 清空默认的页眉页脚

% 设置页眉
\fancyhead[L]{\kaishu \projname}      % 左侧页眉显示项目名称
\fancyhead[C]{\kaishu \reporttitle}    % 中间页眉显示报告标题
\fancyhead[R]{\kaishu \authorname} % 右侧页眉显示作者姓名

% 设置页脚
\fancyfoot[C]{第 \thepage 页，共 \pageref{LastPage} 页} % 中间页脚显示页码

% 去除页眉页脚与正文之间的分隔线
\renewcommand{\headrulewidth}{0.4pt}
\renewcommand{\footrulewidth}{0pt}

% 设置文档背景颜色
\definecolor{mybgcolor}{RGB}{0, 0, 0} % 定义背景颜色
\definecolor{NavyBlue}{RGB}{0,0,128} % 添加颜色定义
% \pagecolor{mybgcolor} % 设置背景颜色

% listing 配置rust
% \lstdefinelanguage{Rust}{
%   keywords={fn, let, mut, struct, enum, impl, pub, use, mod, trait, where},
%   keywordstyle=\color{NavyBlue},
%   comment=[l]{//},
%   commentstyle=\color{gray},
% }
% \lstset{language=Rust}

\begin{document}

% 封面
\begin{titlepage}
    \centering
    

    \includegraphics[width=10cm]{img/SYSULogo.png}

    \vspace{1em}
    %{\Large \college \par}
    \vspace{1em}
    {\Large \kaishu \projname \par}
    \vspace{3em}

    

    
    {\fontsize{40pt}{42pt}\kaishu \selectfont \boldmath \reporttitle\par}
    \vspace*{\fill}

    \begin{center}
    {\Large
    \makebox[5em][s]{学号}:\underline{\makebox[15em][c]{\kaishu \stuno}}\\[1em]
    \makebox[5em][s]{姓名}:\underline{\makebox[15em][c]{\kaishu \authorname}}\\[1em]
    \makebox[5em][s]{专业}:\underline{\makebox[15em][c]{\kaishu \major}}\\[1em]
    \makebox[5em][s]{指导教师}:\underline{\makebox[15em][c]{\kaishu \adviser}}\\[1em]
    \makebox[5em][s]{起始日期}:\underline{\makebox[15em][c]{\kaishu \startdate}}\\[1em]
    \makebox[5em][s]{结束日期}:\underline{\makebox[15em][c]{\kaishu \labenddate}}\\[1em]
    \makebox[5em][s]{实验地点}:\underline{\makebox[15em][c]{\kaishu \labroom}}
    }
    \end{center}

    \vspace*{\fill}
\end{titlepage}

% 目录
\tableofcontents
\newpage

\section{实验概述}

\subsection{实验目的}

\begin{enumerate}
    \item 了解页表的作用、ELF 文件格式、操作系统在 x86 架构的基本启动过程。
    \item 尝试使用 UEFI 加载并跳转到内核执行内核代码。
    \item 实现基于 uart16550 的串口驱动，使用宏启用输出能力、并启用日志系统。
    \item 学习并尝试使用调试器对内核进行调试。\cite{ref1}
\end{enumerate}

\subsection{实验内容}

\begin{enumerate}
    \item 编译内核ELF
    \item 在UEFI中加载内核
    \item UART与日志输出\cite{ref1}
\end{enumerate}

\section{实验步骤与结果}

\subsection{【实验一】配置实验环境}


\section{总结与创新}

\subsection{思考题}


\subsection{加分项}



\subsection{实验总结}

在本次实验中，我配置好了实验环境，学习了基础的Rust，了解了OVMF/UEFI/QEMU等名词，探索了UEFI Shell，编译运行了第一个UEFI程序，并且完成了更加深入的一些加分题。在过程中我对于OS的理解更深一步，也为后续的实验打下了基础。

\section{对实验的改进建议和意见}

\section{附录参考资料}
\begin{thebibliography}{99}  

    \bibitem{ref1}Haoyu Chen. (2024). 中山大学 YatSenOS v2 操作系统实验教程. Website. https://ysos.gzti.me/
    \bibitem{ref2}KIMI AI. (2025). AI协作内容. Website. https://kimi.moonshot.cn/

\end{thebibliography}

\section{附录代码清单}

本实验的代码较多，并且没有很重要的代码，完整代码请参见./Code目录下的代码文件，在对应文件夹下有readme文件说明。

\end{document}